% =============================================================
%  Relativistic modification of Chapter II, §16
%  Cell patterns in relativistic thermal convection
%
%  Original: Chandrasekhar, §16 "The cell patterns"
%  This file extends the analysis to a relativistic fluid
%  using the BDNK causal first-order dissipation framework.
% =============================================================

\section{Cell patterns in relativistic thermal convection}\label{sec:rel-16}

In \S\,\ref{sec:16} it was shown that the theory of marginal stability
at the onset of convection leaves the horizontal cell pattern entirely
unspecified: a disturbance of wave number~$a$ can be decomposed into
orthogonal components in infinitely many ways, and these components may
be superposed with arbitrary amplitudes and phases.  The question we
address here is whether the extension to a relativistic fluid---one in
which the energy density~$\edensity$, pressure~$p$, and inertia
$\enthalpy/c^{2}=(\edensity+p)/c^{2}$ all enter the governing
equations---lifts or preserves this degeneracy at marginal stability,
and how the cell structures are modified when non-linear effects are
taken into account.

%--------------------------------------------------------------
\subsection{Relativistic horizontal eigenvalue problem}\label{sec:rel-16-eigenvalue}
%--------------------------------------------------------------

The separation of the vertical velocity perturbation into horizontal
and vertical factors,
\begin{equation}
\label{eq:rel-2-229}
\delta w = F(x,y)\,W(z),
\end{equation}
carries over to the relativistic problem provided one replaces the
Newtonian horizontal Helmholtz equation
$\nabla_{1}^{2}F = -a^{2}F$
by its relativistic counterpart.  In a static background with metric
$g_{\mu\nu} = \mathrm{diag}(-e^{2\Phi/c^{2}},1,1,1)$, where $\Phi$
is the gravitational potential, the horizontal Laplacian on the
perturbed 3-surface is
\begin{equation}
\label{eq:rel-2-228}
\nabla_{1}^{2}F
  = \left(\frac{\partial^{2}}{\partial x^{2}}
       + \frac{\partial^{2}}{\partial y^{2}}\right)F
  = -a_{\mathrm{rel}}^{2}\,F,
\end{equation}
with the \emph{relativistic horizontal eigenvalue}
\begin{equation}
\label{eq:rel-2-237}
\boxed{\;
a_{\mathrm{rel}}^{2}
  = \bigl(k_{x}^{2}+k_{y}^{2}\bigr)
    \left[1 + \frac{p}{\edensity + p}
             \,\frac{2\Phi}{c^{2}}
         + \order{c^{-4}}\right].
\;}
\end{equation}
Here $k_{x}=2\pi/L_{x}$ and $k_{y}=2\pi/L_{y}$ are the
Newtonian wave numbers in the two horizontal directions.  The
correction factor arises because the spatial part of the perturbed
metric contributes a conformal rescaling of horizontal lengths of
order~$\Phi/c^{2}$.  For a layer of depth~$d$ with
$\Phi\sim gd$, the fractional shift is
$\Delta a^{2}/a^{2}\sim gd/c^{2}$.

\begin{relcorrection}
The classical relation $a^{2}=k_{x}^{2}+k_{y}^{2}$ receives a
multiplicative correction
$1+p\,\Phi/[(\edensity+p)\,c^{2}/2]+\order{c^{-4}}$.  This
correction is \emph{isotropic} in the $(x,y)$-plane and therefore
does not prefer any particular orientation of the wave vector.
\end{relcorrection}

%--------------------------------------------------------------
\subsection{Degeneracy of planforms at marginal stability}\label{sec:rel-16-degeneracy}
%--------------------------------------------------------------

A central result of the classical theory is that the critical Rayleigh
number depends only on the magnitude~$a$ of the horizontal wave
vector, not on its orientation.  This \emph{planform degeneracy}
persists in the relativistic theory for the following reason.

The relativistic linearised perturbation equations, derived from the
conservation laws $\covd T^{\mu\nu}=0$ supplemented by the
BDNK first-order constitutive relations, reduce in the static
background to a system whose coefficients depend on~$z$ alone.
After Fourier decomposition in the horizontal coordinates
$(x,y)\to(k_{x},k_{y})$, every coefficient involves only
$a_{\mathrm{rel}}^{2}=k_{x}^{2}+k_{y}^{2}+\order{c^{-2}}$,
never $k_{x}$ or $k_{y}$ individually.  Consequently the
eigenvalue relation for the relativistic Rayleigh number
\begin{equation}
\label{eq:rel-Ra-marginal}
\Rarel = \Rarel(a_{\mathrm{rel}}^{2})
\end{equation}
is a function of $a_{\mathrm{rel}}^{2}$ alone, and the
minimisation that defines the critical state,
\begin{equation}
\label{eq:rel-Ra-crit}
\Rarel^{(\mathrm{c})} = \min_{a_{\mathrm{rel}}}\Rarel(a_{\mathrm{rel}}^{2}),
\end{equation}
yields a critical eigenvalue $a_{\mathrm{rel,c}}^{2}$ that again
places no restriction on the direction of $\mathbf{k}=(k_{x},k_{y})$.

\begin{causalitycheck}
The BDNK causal frame coefficients enter the vertical eigenvalue
problem but not the horizontal Helmholtz equation.  They therefore
influence $\Rarel^{(\mathrm{c})}$ and $a_{\mathrm{rel,c}}$ but do
not break the horizontal isotropy that underlies the planform
degeneracy.
Causality ($v_{\mathrm{signal}}\le c$) is automatically satisfied
because the BDNK frame coefficients satisfy the coupled
causality inequalities and the perturbation equations are
hyperbolic in the time-like direction.
\end{causalitycheck}

Thus the infinite degeneracy of cell patterns found in the Newtonian
theory---the freedom to superpose plane-wave solutions at arbitrary
orientations and phases so long as
$|\mathbf{k}|=a_{\mathrm{rel,c}}$---\emph{survives intact} in the
relativistic regime.  The cell pattern that is realised must therefore
be determined, as in the classical case, either by symmetry arguments
or by the non-linear dynamics of the finite-amplitude state.

%--------------------------------------------------------------
\subsection{Rolls, rectangles, hexagons, and triangles in the
  relativistic regime}\label{sec:rel-16-patterns}
%--------------------------------------------------------------

Because the horizontal eigenvalue problem~\eqref{eq:rel-2-228} has
exactly the same form as in the Newtonian case---only the value
of~$a$ is shifted---all the cell-pattern solutions of \S\,\ref{sec:16}
carry over with the simple replacement
$a\to a_{\mathrm{rel}}$.  We summarise the relativistic forms for
completeness.

\subsubsection*{(a) Rolls}

\begin{equation}
\label{eq:rel-2-233}
\delta w = W(z)\cos\frac{2\pi}{L_{\mathrm{rel}}}\,x,
\qquad
L_{\mathrm{rel}} = \frac{2\pi}{a_{\mathrm{rel}}}.
\end{equation}
The roll width is $L_{\mathrm{rel}}$, which differs from the
Newtonian value by a fractional amount $\sim gd/c^{2}$.

\subsubsection*{(b) Rectangular and square cells}

\begin{equation}
\label{eq:rel-2-236}
\delta w = W(z)\cos\frac{2\pi}{L_{x}}\,x\;\cos\frac{2\pi}{L_{y}}\,y,
\qquad
a_{\mathrm{rel}}^{2}
  = 4\pi^{2}\!\left(\frac{1}{L_{x}^{2}}+\frac{1}{L_{y}^{2}}\right)
    \relcorr{\Phi/c^{2}}.
\end{equation}
The streamlines in the horizontal plane satisfy the same differential
equation~\eqref{eq:2-246} as in the classical theory and the
cell geometry (sides $L_{x}$, $L_{y}$) is modified only through the
relativistic value of $a_{\mathrm{rel,c}}$.

\subsubsection*{(c) Hexagonal cells}

Christopherson's solution generalises to
\begin{equation}
\label{eq:rel-2-248}
\delta w = \tfrac{1}{3}W(z)\biggl\{
  2\cos\frac{2\pi}{\sqrt{3}\,L_{\mathrm{rel}}}\,x
   \cos\frac{2\pi}{3L_{\mathrm{rel}}}\,y
  + \cos\frac{4\pi}{3L_{\mathrm{rel}}}\,y
\biggr\},
\end{equation}
with $a_{\mathrm{rel}}=4\pi/(3L_{\mathrm{rel}})$.  The hexagonal
symmetry---invariance under $60^{\circ}$ rotation---is preserved because
the relativistic correction to $a^{2}$ is isotropic.  The hexagon side
length $L_{\mathrm{rel}}$ replaces the classical~$L$; the flow topology
(six triangular sub-cells, inscribed hexagon of constant~$w$, etc.)\ is
unchanged.

\subsubsection*{(d) Triangular cells}

As in \S\,\ref{sec:16}, the hexagonal pattern can equally be viewed as
a triangular tessellation with unit-cell side $L_{\mathrm{rel}}\sqrt{3}$.

\subsubsection*{(e) General Bisshopp patterns}

Bisshopp's more general solution~\eqref{eq:2-276} carries over with
the replacement $a\to a_{\mathrm{rel}}$ and the wavelengths
\begin{equation}
\label{eq:rel-2-272}
\lambda_{x}^{\mathrm{rel}}
  = \frac{2\pi}{a_{\mathrm{rel}}}\,\frac{m}{\sqrt{m^{2}-1}},
\qquad
\lambda_{y}^{\mathrm{rel}}
  = \frac{2\pi}{a_{\mathrm{rel}}}\,m.
\end{equation}

%--------------------------------------------------------------
\subsection{Non-linear effects: relativistic corrections to
  amplitude equations}\label{sec:rel-16-nonlinear}
%--------------------------------------------------------------

Although the linear theory does not select the cell pattern, the
weakly non-linear theory does---through the \emph{amplitude equations}
(Landau or Ginzburg--Landau type) that govern the slow evolution of
finite-amplitude perturbations near the critical Rayleigh number.

Denote by $A_{j}$ ($j=1,\ldots,N$) the complex amplitudes of~$N$
plane-wave modes whose wave vectors $\mathbf{k}_{j}$ all have
$|\mathbf{k}_{j}|=a_{\mathrm{rel,c}}$.  To cubic order the amplitude
equations take the standard form
\begin{equation}
\label{eq:rel-amplitude}
\tau_{\mathrm{rel}}\,\frac{\partial A_{j}}{\partial t}
  = \sigma\,A_{j}
  - \ell_{j}\,|A_{j}|^{2}A_{j}
  - \sum_{k\neq j} g_{jk}\,|A_{k}|^{2}A_{j}
  + h_{j}\,A_{m}^{*}A_{n}^{*}\,\delta_{\mathbf{k}_{j}+\mathbf{k}_{m}+\mathbf{k}_{n},\,0},
\end{equation}
where $\sigma\propto(\Rarel-\Rarel^{(\mathrm{c})})$ is the
supercriticality parameter.

\paragraph{Relativistic modifications of the Landau coefficients.}
The self-interaction coefficient $\ell_{j}$, the cross-coupling
$g_{jk}$, and the resonant triad coefficient $h_{j}$ all receive
$\order{v^{2}/c^{2}}$ corrections compared with their Newtonian
values.  These corrections enter through:
\begin{enumerate}
\item \emph{Relativistic inertia}: the role of the rest-mass
  density~$\rdensity$ in the non-linear advection term
  $u^{\alpha}\covda u^{\mu}$ is replaced by the enthalpy
  $\enthalpy/c^{2}=(\edensity+p)/c^{2}$, enhancing the effective
  inertia by a factor $1+p/(\edensity c^{2})+\order{c^{-4}}$.
\item \emph{Causal heat transport}: the BDNK frame coefficients
  introduce frequency-dependent (dispersive) corrections to the
  effective thermal diffusivity that modify the saturation
  coefficients $\ell_{j}$ and $g_{jk}$.
\item \emph{Compressibility}: relativistic fluids with
  $c_{\mathrm{s}}/c$ not negligible develop a coupling between
  the thermal and acoustic modes that has no Newtonian counterpart.
  This generates an additional cubic term of
  order~$(c_{\mathrm{s}}/c)^{2}$ in~\eqref{eq:rel-amplitude}.
\end{enumerate}

\paragraph{Pattern selection.}
The classical result of Busse and others---that rolls are the stable
planform for Boussinesq convection between stress-free boundaries
with up/down symmetry---carries over to the relativistic case as
long as the Boussinesq-type symmetry is preserved.  However, two
mechanisms can break this symmetry relativistically:
\begin{itemize}
\item A non-negligible gravitational redshift across the layer
  ($gd/c^{2}\not\ll 1$) introduces a preferred vertical direction
  in the energetics, making $h_{j}\neq 0$ and favouring hexagonal
  over roll-like patterns.
\item Bulk-viscous effects ($\bulkP\neq 0$) associated with the
  equation of state of dense matter couple the mean dilatation to
  the convective amplitude, again generating a hexagonal tendency.
\end{itemize}

%--------------------------------------------------------------
\subsection{Physical discussion and astrophysical
  relevance}\label{sec:rel-16-discussion}
%--------------------------------------------------------------

\subsubsection*{Laboratory fluids}

For terrestrial experiments the relativistic corrections are
negligible.  Taking water as an example:
\begin{equation}
\label{eq:rel-lab-estimate}
\frac{gd}{c^{2}}
  \sim \frac{10\;\mathrm{m\,s^{-2}}\times 10^{-2}\;\mathrm{m}}
       {(3\times 10^{8}\;\mathrm{m\,s^{-1}})^{2}}
  \sim 10^{-18}.
\end{equation}
The fractional correction to $a^{2}$ and to the Landau coefficients
is of this order---far beyond any conceivable measurement.  The
classical cell-pattern analysis of \S\,\ref{sec:16} therefore
applies without modification to all laboratory B\'enard experiments.

\subsubsection*{Neutron-star oceans}

In the outer envelope of a neutron star, $g\sim 10^{14}\;\mathrm{m\,s^{-2}}$
and $d\sim 1$--$10\;\mathrm{m}$, giving
\begin{equation}
\label{eq:rel-ns-estimate}
\frac{gd}{c^{2}} \sim 10^{-3}\text{--}10^{-2}.
\end{equation}
Corrections at the percent level can affect the critical wave number
and thus the cell size of convective patterns in accreting neutron-star
oceans.  Moreover, the equation of state of dense matter yields
$c_{\mathrm{s}}^{2}/c^{2}\sim 0.1$--$0.3$, so the compressibility
coupling in~\eqref{eq:rel-amplitude} is no longer negligible and may
shift the non-linear pattern selection toward hexagons.

\subsubsection*{Quark--gluon plasma}

In the quark--gluon plasma produced in heavy-ion collisions,
$c_{\mathrm{s}}^{2}=c^{2}/3$ (conformal limit) and the fluid is
ultra-relativistic.  While the geometry is far from a classical
B\'enard layer, the mathematical structure of the amplitude
equations~\eqref{eq:rel-amplitude} applies to any convective
instability in a relativistic fluid.  The Landau coefficients
receive $\order{1}$ corrections; the Newtonian expressions are
not even a useful starting approximation.

\subsubsection*{Early-universe plasma}

During the radiation-dominated era, the cosmic fluid has
$p=\edensity/3$ and $\enthalpy/c^{2}=4\edensity/(3c^{2})$.
Convective instabilities driven by compositional gradients (e.g.\
lepton-number inhomogeneities) would obey the full relativistic
amplitude equations with $\order{1}$ relativistic effects.  The
planform degeneracy still holds at linear order, but the pattern
selected by the non-linear dynamics is governed entirely by the
relativistic coefficients.

%--------------------------------------------------------------
\subsection{Summary}\label{sec:rel-16-summary}
%--------------------------------------------------------------

\begin{enumerate}
\item The horizontal eigenvalue relation
  $a^{2}=k_{x}^{2}+k_{y}^{2}$ receives a multiplicative
  relativistic correction~\eqref{eq:rel-2-237} that is isotropic
  and of order $gd/c^{2}$.
\item The \emph{degeneracy of planforms at marginal stability
  persists} in the relativistic theory: the critical Rayleigh number
  depends only on $a_{\mathrm{rel}}^{2}$, not on the direction
  of~$\mathbf{k}$.
\item All classical cell-pattern solutions (rolls, rectangles,
  hexagons, triangles, Bisshopp patterns) carry over with the
  replacement $a\to a_{\mathrm{rel}}$.
\item Non-linear amplitude equations receive $\order{v^{2}/c^{2}}$
  corrections from relativistic inertia, causal heat transport,
  and compressibility.  These corrections can shift the preferred
  pattern from rolls to hexagons when $gd/c^{2}$ or
  $c_{\mathrm{s}}^{2}/c^{2}$ is not negligible.
\item For laboratory fluids the corrections are unmeasurably small
  ($\sim 10^{-18}$); for neutron-star oceans they reach
  $\sim 10^{-2}$; for quark--gluon plasma and the early universe
  they are of order unity.
\end{enumerate}

%--------------------------------------------------------------
\subsection{Astrophysical applications: quantitative pattern selection}
\label{sec:rel-16-applications}
%--------------------------------------------------------------

\paragraph{Neutron star ocean.}
With $g\approx 2.4\times 10^{14}\;\mathrm{cm\,s^{-2}}$ and $d\approx 10$\,m:
$gd/c^{2}\approx 2.7\times 10^{-4}$, cell size correction $\sim 0.015\%$.
However, $c_s^{2}/c^{2}\approx 0.15$ generates non-negligible
compressibility coupling in the amplitude equation.

\paragraph{QGP fireball.}
With $c_s^{2}=c^{2}/3$ and $\xi\approx 1/3$: self-interaction Landau
coefficient enhanced by $4/3$, cross-coupling by $3/2$.  Compressibility
coupling $h\sim 1/3$ strongly favours hexagonal patterns.

\paragraph{Early universe.}
$gd/c^{2}\sim 0.1$ for sub-horizon perturbations; classical Busse
theory is inapplicable.

\begin{figure}[htbp]
\centering
\includegraphics[width=0.95\textwidth]{fig_cell_pattern_selection.pdf}
\caption{(a)~Pattern selection phase diagram: rolls stable at low
$gd/c^{2}$ and $c_s^{2}/c^{2}$; hexagons preferred at larger values.
(b)~Landau coefficients vs.\ $\xi$ at $gd/c^{2}=0.01$.}
\label{fig:cell-pattern-selection}
\end{figure}

\section*{References}

\begin{enumerate}
\item S.~Chandrasekhar, \textit{Hydrodynamic and Hydromagnetic Stability}, OUP (1961).
\item F.~H.~Busse, \textit{Rep.\ Prog.\ Phys.}\ \textbf{41}, 1929 (1978).
\item F.~S.~Bemfica et al., \textit{Phys.\ Rev.\ D} \textbf{98}, 104064 (2018).
\item P.~Kovtun, \textit{JHEP} \textbf{10}, 034 (2019).
\item P.~Haensel, A.~Y.~Potekhin, and D.~G.~Yakovlev,
  \textit{Neutron Stars 1}, Springer (2007).
\item M.~G.~Alford, S.~Mahmoodifar, and K.~Schwenzer,
  \textit{Phys.\ Rev.\ D} \textbf{85}, 044051 (2012).
\item P.~Romatschke and U.~Romatschke, \textit{Relativistic Fluid Dynamics}, CUP (2019).
\item L.~Rezzolla and O.~Zanotti, \textit{Relativistic Hydrodynamics}, OUP (2013).
\end{enumerate}
